\section{Introduction and problem}
\label{sec:intro}

When a driving policy fails in closed loop, the team that owns it has to decide
what data to collect next. To do that it needs the failures sorted into groups
that share a cause: a group of collisions with pedestrians walking in the lane, a
group of stalls in front of a fully blocked road, and so on. The usual way to get
those groups is to write a list of candidate factors (weather, junction type,
object class) and measure how much the policy degrades along each. That only finds
failures someone already thought to write down.

This study asks a narrower question that can be measured. If we take the policy's
internal activations on its own driven routes, and the pass/fail outcome of each
route, can a method with no scenario label sort the failing routes into the
groups that a benchmark already knows about? We call this \emph{label-free
failure discovery}. A method sees activations, the outcome of each route, and
signals a fleet logs (speed, infractions). It never sees the scenario name, the
route configuration, or any field that encodes the failure type.

\paragraph{Why a benchmark with known groups.}
We cannot score discovery of failures nobody knows. So we score recovery of
failure groups that are known to us and hidden from every method, and we read
that as a proxy for finding unknown ones. Fail2Drive~\cite{gerstenecker2026fail2drive}
suits this because each of its routes comes as a pair: a base route in the
training distribution and a twin on the same road that adds one unseen scenario.
% published: 100 route pairs, 17 scenario types (CITATION_AUDIT.md, fail2drive row)
It has 100 such pairs over 17 scenario types. A scenario type whose twins fail
much more often than their bases is a \emph{gap}: a failure group with a known
cause. The proxy is stated here because a reader may otherwise take it for the
real goal.

\paragraph{The comparison that matters.}
A discovery method is only useful if it beats what a fleet already has. A fleet
logs the car's speed and the infractions it incurred, so our main rival groups
failures by those signals alone (B3-behaviour, \S\ref{sec:rivals}). An earlier
rival that also read the simulator's exact distances to nearby objects was
retired, because no deployed car has those distances. The question is therefore
whether activations add anything over behaviour, not whether they beat chance.

\paragraph{Models.}
We develop methods on SimLingo~\cite{Renz2025cvpr}, a vision-language-action
driving model. We freeze the finalists and score them on TransFuser++
\cite{jaeger2023hidden,zimmerlin2024hidden}, a camera and LiDAR model, and on
UniAD, a modular camera-only model. UniAD's outcomes are blind to every method
until the pre-registered read (\S\ref{sec:prereg}). All driving runs in CARLA
\cite{carla2017}.

\paragraph{What this draft contains.}
Section~\ref{sec:data} describes the data pools and their roles.
Section~\ref{sec:methods} describes the search framework, the two method families,
the rivals and the scoring rules at a level a reader can reimplement.
Section~\ref{sec:prereg} gives the pre-registration history, including every
post hoc amendment. Section~\ref{sec:whatbroke} lists what went wrong that a
reader needs in order to trust or doubt the numbers. Section~\ref{sec:related}
places the work. There is no results section yet.

\paragraph{What this study does not claim.}
It makes no claim about real vehicles, about transfer to other domains, or that
driving is a new setting for failure analysis. It tests recovery of known groups,
not discovery of unknown ones. Every figure in the inherited planning documents
that was a target rather than a measurement is left out.
